\subsection{窄收敛的紧性判据}

定义 2. --- 设 T 是拓扑空间，H 是 \(\mathcal{M}^{\mathrm{b}}(\mathrm{T})\) 的子集；称 H 满足 Prokhorov 条件，如果

\begin{enumerate}
\def\labelenumi{\alph{enumi})}
\item
  \(\sup_{\mu \in \mathrm{H}}|\mu |(1) < +\infty ;\)
\item
  对每个数 \(\varepsilon > 0\)，存在紧子集 \(\mathbf{K}_{\varepsilon}\)（属于 \(\mathbf{T}\)），使得
\end{enumerate}

\[
| \mu | (\mathrm{T} - \mathrm{K} _ {\varepsilon}) \leqslant \varepsilon \quad \text { for every measure } \mu \in \mathrm{H}.\tag{6}
\]

可以证明，若 T 完全正则，则条件 a) 和 b) 等价于以下条件：存在 T 上满足 \(f \geqslant 1\) 的实函数 f，使得满足 \(f(t) \leqslant c\) 的 T 中点 t 所成的集合对每个 \(c \in R_{+}\) 都是紧的（这特别蕴含 f 下半连续），并且 \(\sup_{\mu \in \mathbf{H}} |\mu|(f) < +\infty\)。此外，当 \(\mathbf{T}\) 局部紧时，要求 \(f\) 连续可得到一个等价的陈述（参见习题 10）。

命题 11. --- 设 T 是完全正则空间，H 是 \(\mathcal{M}^{b}(T)\) 的满足 Prokhorov 条件的子集；则其闭包 \(\overline{H}\) 在 \(\mathcal{M}^{b}(T)\) 中也满足 Prokhorov 条件。

事实上，由第 3 号命题 6，映射 \(\mu \mapsto |\mu|^{\bullet}(1)\)、\(\mu \mapsto |\mu|^{\bullet}(\mathrm{T} - \mathrm{K}_{\varepsilon})\) 在 \(\mathcal{M}^b (\mathrm{T})\) 上下半连续。

Prokhorov 条件的重要性来自以下定理；其逆命题将在稍后研究（定理 2）。

定理 1（Prokhorov）。 --- 设 T 是完全正则空间，H 是满足 Prokhorov 条件的 \(\mathcal{M}^{b}(\mathrm{T})\) 的子集；则 H 在 \(\mathcal{M}^{b}(\mathrm{T})\) 中关于窄拓扑相对紧。

不妨设 T 是紧空间 X 的子空间；令 i 为从 T 到 X 的典范嵌入。另一方面，由命题 11，可设 H 在 \(\mathcal{M}^{b}(T)\) 中是闭的。于是只须证明 H 上的每个超滤子都在 \(\mathcal{M}^{b}(T)\) 中收敛。

先处理 \(\mathrm{H}\subset\mathcal{M}_{+}^{b}(\mathrm{T})\) 的情形。由假设，\(\mu\in H\) 的测度总质量有界，故 \(i(\mu)\) 关于 U 在 \(\mathcal{M}_{+}(\mathrm{X})\) 中淡收敛到某个测度 \(\nu\in\mathcal{M}_{+}(\mathrm{X})\)（第 III 章 §1, No.~9, 命题 15 的推论 2）；由第 3 号命题 8，问题归结为证明 \(\nu\) 集中于 T。现在取数 \(\varepsilon\) \textgreater0，并取满足公式 (6) 的 T 的紧子集 \(K_{\varepsilon}\)。由于 \(X-K_{\varepsilon}\) 是 X 中的开集，由第 3 号命题 6 在 X 中的应用，得到不等式

\[
\begin{array}{r l} \nu^ {\bullet} (\mathrm{X} - \mathrm{T}) & \leqslant \nu^ {\bullet} (\mathrm{X} - \mathrm{K} _ {\varepsilon}) \leqslant \liminf _ {\mu , \mathfrak {U}} i (\mu) ^ {\bullet} (\mathrm{X} - \mathrm{K} _ {\varepsilon}) \\ & = \liminf _ {\mu , \mathfrak {U}} \mu^ {\bullet} (\mathrm{T} - \mathrm{K} _ {\varepsilon}) \leqslant \varepsilon ; \end{array}
\]

由于 \(\varepsilon > 0\) 任意，定理在这一特殊情形下得证。

现在处理一般情形；对 T 上的每个测度 \(\mu\)，置

\[
a _ {1} (\mu) = \mathcal {R} (\mu) ^ {+}, a _ {2} (\mu) = \mathcal {R} (\mu) ^ {-}, a _ {3} (\mu) = \mathcal {I} (\mu) ^ {+}, a _ {4} (\mu) = \mathcal {I} (\mu) ^ {-};
\]

由于 \(\mu = a_{1}(\mu) - a_{2}(\mu) + ia_{3}(\mu) - ia_{4}(\mu)\)，只须证明映射 \(a_{j}\)（\(j = 1,2,3,4\)）关于 \(\mathfrak{U}\) 窄收敛。但是，由测度 \(\mathbf{H}_j\)（其中 \(a_{j}(\mu)\) 遍历 \(\mu\)）所成的集合 \(\mathbf{H}\)，根据关系 \(|a_j(\mu)| \leqslant |\mu|\) 满足 Prokhorov 条件，且包含于 \(\mathcal{M}_+^b(\mathrm{T})\)；因此根据特殊情形，它在 \(\mathcal{M}_+^b(\mathrm{T})\) 中相对紧，于是定理立即得出。

推论。 --- 设 T 是完全正则空间 X 的子空间，H 是满足 Prokhorov 条件的 \(\mathcal{M}^{b}(T)\) 的子集。若 i 表示从 T 到 X 的典范嵌入，则映射 \(\mu\mapsto i(\mu)\) 从 \(\mathcal{M}^{b}(T)\) 到 \(\mathcal{M}^{b}(X)\) 在 H 上的限制是 H 到其像的同胚。

只须处理 H 闭的情形（命题 11），此时 H 紧；结论由映射 \(\mu \mapsto i(\mu)\) 连续且单射这一事实立即得到。

回忆一下，这一结果对 \(\mathcal{M}_{+}^{b}(\mathrm{T})\) 的任意子集也成立（第 3 号命题 8）。

定理 2. --- 设 T 是局部紧空间或波兰空间，H 是 \(\mathcal{M}_{+}^{b}(T)\) 的相对紧子集；则 H 满足 Prokhorov 条件。

不妨设 H 闭，从而紧。测度 \(\mu \in H\) 的总质量显然有界，因为映射 \(\mu \mapsto \mu(1)\) 连续，问题归结为证明定义 2 的断言 b)。

先设 T 局部紧。取数 \(\varepsilon\) \textgreater0。对每个测度 \(\mu\in H\)，取 T 中的紧集 \(K_{\mu}\)，使得 \(\mu^{\bullet}(T-K_{\varepsilon})<\varepsilon\)，再取 \(U_{\mu}\) 的相对紧开邻域 \(K_{\mu}\)。由于映射 \(\lambda\mapsto\lambda^{\bullet}(T-U_{\mu})\) 在 \(\mathcal{M}_{+}^{b}(T)\) 上上半连续（第 3 号命题 6），所以由满足 \(V^{\mu}\) 的 H 中测度 \(\lambda\in H\) 所成的集合 \(\lambda^{\bullet}(T-U_{\mu})<\varepsilon\) 是 H 中 \(\mu\) 的邻域。因此存在 H 的有限子集 \(H'\)，使得集合 \(V^{\mu}\)（\(\mu\in H'\)）覆盖 H。记 \(\bigcup_{\mu\in H'}\overline{U}_{\mu}\) 为紧集，则对所有 \(\lambda^{\bullet}(T-K)<\varepsilon\) 有 \(\lambda\in H\)。

再设 T 是波兰空间。可不失一般性地假设 T 是紧空间 X 中递减开集序列 \((\mathrm{T}_{p})_{p\geqslant1}\) 的交（GT, IX, §6, No.~1, 定理 1 的推论 1）。令 \(i_{p}\) 为从 T 到 \(T_{p}\) 的嵌入，令 \(H_{p}\) 为形如 \(i_{p}(\lambda)\)（其中 \(\lambda\in H\)）的测度所成的集合；由于 \(H_{p}\) 在 \(\mathcal{M}_{+}^{b}(T_{p})\) 中紧，应用前面对局部紧空间 \(K_{p}\subset T_{p}\) 的结果可知，存在紧集 \(\nu^{\bullet}(T_{p}-K_{p})\leqslant\varepsilon2^{-p}\)，使得 \(\nu\in H_{p}\) 对每个测度 \(T_{p}\) 成立。因此也有 \(\nu^{\bullet}\big(\mathrm{T}-(\mathrm{T}\cap\mathrm{K}_{p})\big)\leqslant\varepsilon2^{-p}\)，最后 \(\lambda^{\bullet}\big(\mathrm{T}-(\mathrm{T}\cap\mathrm{K}_{p})\big)\leqslant\varepsilon2^{-p}\) 对每个测度 \(\lambda\in H\) 成立。现在置 \(K=\bigcap_{p}K_{p}\)；集合 K 紧且包含于 T，并且 \(\lambda\in H\) 对每个测度 \(\lambda^{\bullet}(T-K)\leqslant\sum_{p}\lambda^{\bullet}\big(\mathrm{T}-(\mathrm{T}\cap\mathrm{K}_{p})\big)\leqslant\sum_{p}\varepsilon2^{-p}=\varepsilon\) 成立。因此 Prokhorov 条件得到验证。

